\section*{5. 指示薬と水溶液の色}
\begin{enumerate}
    \item 次の表の空欄に、水溶液の性質によって変化する指示薬の色をそれぞれ答えなさい。\\
    （※色が変化しない場合も、そのときの色を答えること）
\end{enumerate}
\begin{center}
\renewcommand{\arraystretch}{1.8}
\begin{tabular}{|c|c|c|c|}
\hline
指示薬 & 酸性 & 中性 & アルカリ性 \\ \hline
青色リトマス紙 & \hspace{3cm} &  \hspace{3cm}  &  \hspace{3cm}  \\ \hline
赤色リトマス紙 &  \hspace{3cm}  &  \hspace{3cm}  &  \hspace{3cm}  \\ \hline
BTB溶液 &  \hspace{3cm} &  \hspace{3cm}  & \hspace{3cm} \\ \hline
フェノールフタレイン溶液 &  \hspace{3cm}  &  \hspace{3cm} &  \hspace{3cm}  \\ \hline
pH試験紙 &  \hspace{3cm}  &  \hspace{3cm}  & \hspace{3cm}  \\ \hline
\end{tabular}
\end{center}
